% ProbabilityStatistics_A_#11
%%%%%%%%%%%%%%%

\section{法則収束と特性関数(Fourier変換)}
\subsection{確率測度の空間}
分布と特性関数
法則収束と分布関数

%\begin{soudef}
[法則収束]
    確率変数の列$\{X_n\}$と確率変数$X$に対して
\begin{equation}
    \lim_{n\to \infty} \mathbb{E}[F(X_n)]= \mathbb{E}[F(X)],\ \ \forall F\in C_b(\R),\ \ \mathrm{有界連続}
\end{equation}
    をみたすとき，$X_n$は$X$に法則収束(in lawで収束する)といい，
\begin{equation}
    \lim_{n\to \infty} X_n= X\ \ \mathrm{in law}
\end{equation}
    で表す．
%\end{soudef}
%\begin{souremark}
    もし$X_n$が$X$に概収束するならば，法則収束する．\\
    確率収束しているときも，法則収束がいえる．
%\end{souremark}
    これは次のことから従う．
%\begin{souprop}
    $\{a_n\}_{n= 1}^\infty \subset \R$は次を満たすとする．\\
    ある$\alpha\in \R$に対して任意の部分列$\{\alpha_{n(k)}\}_{k= 1}^\infty$に対してさらに部分列$\{a_{n'(k)}\}_{k= 1}^\infty$をとれば収束し，収束先は$\alpha$である．\\
    このとき$\lim_{n\to \infty} a_n= \alpha$
%\end{souprop}

%\begin{souprop}
    $(\Omega, \F, \mathcal{P})$上の確率変数の列$\{X_n\}$と確率変数$X$について以下は同値
\begin{enumerate}
    \item 
    $X_n$は$X$に法則収束
    \item\label{step:x= aで分布関数が収束} 
    $\forall a\in \R$に対し，もし
\begin{equation}
    \P(X= a)= 0
\end{equation}
    ならば
\begin{equation}
    \lim_{n\to \infty} \mathcal{P}(X_n\le a)= \mathcal{P}(X\le a)
\end{equation}
\end{enumerate}
%\end{souprop}

%\begin{soucor}
    $X$の分布が連続型ならば，\eqref{step:x= aで分布関数が収束}の代わりに\eqref{step:「x= aで分布関数が収束」の代用}でよい．
\begin{enumerate}
    \item [$(2)'$]\label{step:「x= aで分布関数が収束」の代用}
    $\forall a\in \R$に対し，
\begin{equation}
    \lim_{n\to \infty} \mathcal{P}(X_n\le a)= \mathcal{P}(X\le a)
\end{equation}
\end{enumerate}
%\end{soucor}

\subsection{特性関数}
%\begin{soudef}
[特性関数]
    $\mu$を$(\R^n, \mathcal{B(\R^n)})$上の確率測度とする．このとき
\begin{equation}
    \phi_\mu(t)= \int_{\R^n} \exp(i(t, x))\mu(dx)
\end{equation}
    を$\mu$の特性関数という．
%\begin{souremark}
    実数値確率変数$X$に対して，その分布を$\mu$とすると
\begin{equation}
    \phi_\mu(t)= \mathbb{E}[\exp(i(t, x))]
\end{equation}
    となる．これを$X$の特性関数という．
%\end{souremark}
%\begin{souremark}
    複素数値の積分は実部・虚部を分けて考える．\\
    $|\exp(i(t, x))|= 1,\ \ x, y\in \R$
%\end{souremark}
%\end{soudef}

%\begin{souprop}
    実数値確率変数$X$は
\begin{equation}
    \mathbb{E}[|X|^n]< \infty
\end{equation}
    を満たすとする．このとき$\phi_\mu$は$n$回連続微分で 
\begin{equation}
    \phi^{(n)}_\mu(t)= \mathbb{E}[(iX)^n\exp(itX)],\ \ t\in \R
\end{equation}
    特に
\begin{equation}
    \phi^{(n)}_\mu(0)= i^n\mathbb{E}[X^n]
\end{equation}
    特性関数は分布$\mu$の情報を多分に含む．
%\end{souprop}
%\begin{souprop}
    $X$がパラメータ$\lambda$のPoisson分布に従う．つまり
\begin{equation}
    \mathcal{P}(X= k)= e^{-\lambda}\frac{\lambda^k}{k!},\ \ k\ge 0
\end{equation}
    とする．このとき$X$の特性関数$\phi_X$は
\begin{equation}
    \phi_X(t)= \exp(\lambda(e^{it}- 1))
\end{equation}
%\end{souprop}
\begin{proof}
\begin{align}
    \phi_X(t)
    &= \mathbb{E}[\exp(itX)]\\
    &= \sum_{k= 0}^\infty e^{itk}e^{-\lambda}\frac{\lambda^k}{k!}\\
    &= e^{-\lambda}\sum_{k= 0}^\infty \frac{1}{k!}(\lambda e^{it})^k\\
    &= e^{-\lambda}\cdot e^{\lambda e^{it}}\\
    &= (右辺)
\end{align}
\end{proof}

%\begin{souprop}
    確率変数$X, Y$は独立でそれぞれの特性関数を$\phi_X, \phi_Y$とする．このとき$X+ Y$の特性関数$\phi_{X+ Y}$は
\begin{equation}
    \phi_{X+ Y}= \phi_X(t)\cdot \phi_Y(t)
\end{equation}
%\end{souprop}
%\begin{souprop} \label{thm:aX+bの特性関数}
    確率変数$X$の特性関数を$\phi_X$とする．$a, b\in \R$に対して，$aX+ b$の特性関数$\phi_{aX+ b}$は
\begin{equation}
    \phi_{aX+ b}(t)= e^{itb} \phi_X(at)
\end{equation}
%\end{souprop}

%\begin{souprop}
    $m\in \R, v> 0$とする．$X\sim \mathcal{N}(m, v)$のとき
\begin{equation}
    \phi_X(t)= \exp{\left(imt- \frac{1}{2}vt^2\right)}
\end{equation}
%\end{souprop}
\begin{proof}
    まず$m= 0, v= 1$のときを示す．このとき
\begin{equation}
    \phi_X(t)= \mathbb{E}[e^{itX}]= \mathbb{E}[\cos{tX}]+ i\mathbb{E}[\sin{tX}]
\end{equation}
    ここで
\begin{equation}
    \mathbb{E}[\sin{tX}]= \int_{-\infty}^\infty \sin{tX}\cdot \frac{1}{\sqrt{2\pi}}e^{-\frac{1}{2}x^2}dx= 0
\end{equation}
    また，
\begin{equation}
    \phi'_X(t)= \mathbb{E}[(iX)e^{itX}]= -\mathbb{E}[X\sin{tX}]
\end{equation}
    であって
\begin{align}
    \mathbb{E}[X\sin{tX}]
    &= \int_{-\infty}^\infty x\sin{tx}\frac{1}{\sqrt{2\pi}}e^{-\frac{1}{2}x^2}dx\\
    &= -\int_{-\infty}^\infty \sin{tx}\left(\frac{1}{\sqrt{2\pi}}e^{-\frac{1}{2}x^2}\right)'dx\\
    &= \int_{-\infty}^\infty t\cos{tx}\frac{1}{\sqrt{2\pi}}e^{-\frac{1}{2}x^2}dx\\
    &= t\phi_X(t)\\
    \therefore \phi'_X(t)&= -t\phi_X(t)
\end{align}
    $\phi_X(0)\ne 0$のとき
\begin{equation}
    \frac{\phi'(t)}{\phi_X(t)}= -t
    \iff \log|\phi_X(t)|= -\frac{1}{2}t^2+ C
    \iff \phi_X(t)= \exp{\left(-\frac{1}{2}t^2+ C\right)}
\end{equation}
    ここで$\phi_X(0)= 1$より$C= 0$\\
    $\therefore \phi_X(t)= \exp{\left(-\frac{1}{2}t^2\right)}$
    あとは$X\sim \mathcal{N}(m, v)$のとき
\begin{equation}
    z=: \frac{X- m}{\sqrt{v}}
\end{equation}
    とおくと$z\sim \mathcal{N}(0, 1)$で
\begin{equation}
    X\sim \sqrt{v}z+ m
\end{equation}
    となる．定理\ref{thm:aX+bの特性関数}から結論を得る．
\end{proof}